%! AP Statistics — Unit 4, Lesson 4.2 — Student Packet Cover Sheet
\documentclass[10pt]{article}
\usepackage{apstats-article}
\usepackage{apstats-boxes}
\usepackage{ltablex}
\keepXColumns

\begin{document}

% Full-bleed navy banner
\begin{tikzpicture}[remember picture, overlay]
  \fill[navy] (current page.north west)
    rectangle ([yshift=-0.9in]current page.north east);
\end{tikzpicture}
\vspace{-0.8in}

\begin{center}
  {\color{white}\textbf{\LARGE AP Statistics}}\\[4pt]
  {\color{skymid}\large\textbf{Unit 4: Probability, Random Variables, and Probability Distributions}}\\[6pt]
  {\color{white}\textbf{\large Lesson 4.2 \quad Estimating Probabilities Using Simulation}}
\end{center}

\vspace{0.22in}
\namedateperiod
\vspace{0.18in}

\begin{learningtargetbox}
\small
\begin{enumerate}[leftmargin=1.5em, itemsep=5pt]
  \item \ldots{} design and run a simulation of a random process using an appropriate chance device.
  \item \ldots{} use relative frequency of simulated outcomes to estimate probability of an event.
  \item \ldots{} explain why more trials produce a more accurate probability estimate (Law of Large Numbers).
\end{enumerate}
\end{learningtargetbox}

\vspace{0.14in}

\begin{tocbox}
\small
\renewcommand{\arraystretch}{1.5}
\begin{tabularx}{\linewidth}{@{} c l X r @{}}
\rowcolor{sky}
\textbf{\#} & \textbf{Component} & \textbf{Description} & \textbf{Score} \\
\midrule
1 & Warm-Up       & Spiral review + relative frequency preview             & \blank{1.2cm} \\
2 & Guided Notes  & Simulation design; law of large numbers                & \blank{1.2cm} \\
3 & Group Activity & Design and run a simulation (genetics scenario)        & \blank{1.2cm} \\
4 & Exit Ticket   & Marble-drawing simulation analysis                     & \blank{1.2cm} \\
5 & Homework      & Full simulation design; law of large numbers reflection & \blank{1.2cm} \\
\midrule
  & \multicolumn{2}{r}{\textbf{Total}} & \blank{1.2cm} \\
\end{tabularx}
\end{tocbox}

\vspace{0.14in}

\begin{remindbox}
\small
The \textbf{most important idea in Lesson 4.2}: simulation is a tool for estimating
probability --- the more trials you run, the closer your estimate gets to the true value.

\vspace{6pt}
\begin{tabularx}{\linewidth}{@{} >{\centering\arraybackslash\columncolor{goldbg}}p{0.06\linewidth}
                                    X @{}}
\textbf{\large!} &
\textbf{Simulation = modeling randomness.}\enspace
You do not need a formula to estimate a probability. Run the random process many times
and compute the relative frequency of the outcome you care about. \\[6pt]
\textbf{\large!} &
\textbf{Law of Large Numbers.}\enspace
The simulated (empirical) probability will fluctuate early on but will eventually
stabilize near the true probability as trials increase. \\[6pt]
\textbf{\large!} &
\textbf{Design matters.}\enspace
A good simulation correctly maps each possible outcome of the real process to a chance
device value, and accounts for all possible outcomes. \\
\end{tabularx}
\end{remindbox}

\end{document}
